\documentclass[11pt,a4paper]{article}

% ===================== PACKAGES =====================
\usepackage[utf8]{inputenc}
\usepackage[T1]{fontenc}
\usepackage[provide=*,french]{babel}

% Guillemets fran\c{c}ais (repli si babel-french n'est pas install\'e)
\providecommand{\og}{\leavevmode\guillemotleft{}~}
\providecommand{\fg}{\unskip~\guillemotright{}}
\usepackage[margin=2.5cm]{geometry}
\usepackage{graphicx}
\usepackage{float}
\usepackage{listings}
\usepackage{hyperref}
\usepackage{enumitem}
\usepackage{titling}
\usepackage[font=small,labelfont=bf,labelsep=endash]{caption}
\usepackage{amssymb}
\usepackage{xcolor}
\definecolor{blueESP}{HTML}{0040FF}

\graphicspath{{images/}{./}}

\hypersetup{
    hidelinks
}

% ===================== CONFIG LISTINGS (sans couleur) =====================
\lstset{
    basicstyle=\ttfamily\small,
    breaklines=true,
    frame=single,
    columns=fullflexible,
    showstringspaces=false,
    tabsize=2
}

% ===================== PAGE DE GARDE =====================
\title{Rapport de TP 1 \\ Découverte et administration d'un système Linux}
\author{}
\date{}

\begin{document}

\begin{titlepage}
\centering

%%% Bandeau institutionnel
\noindent
\begin{minipage}[c]{0.17\textwidth}
    \centering
    \includegraphics[height=1.5cm]{logoESPA.jpg}
\end{minipage}
\hfill
\begin{minipage}[c]{0.64\textwidth}
    \centering
    {\bfseries REPOBLIKAN'I MADAGASIKARA}\\[2pt]
    {\scriptsize\bfseries FITIAVANA - TANINDRAZANA - FANDROSOANA}\\[5pt]
    {\scriptsize\bfseries MINIST\`ERE DE L'ENSEIGNEMENT SUP\'ERIEUR ET DE\\[1pt] LA RECHERCHE SCIENTIFIQUE}
\end{minipage}
\hfill
\begin{minipage}[c]{0.17\textwidth}
    \centering
    \includegraphics[height=1.5cm]{logoUNA.jpg}
\end{minipage}

\vspace{1.4cm}

{\bfseries \'ECOLE SUP\'ERIEURE POLYTECHNIQUE D'ANTSIRANANA}

\vspace{1.2cm}

{\bfseries Mention Master STIC}

\vspace{1cm}

%%% Bandeau du type de document
\noindent\rule{\textwidth}{1.6pt}\\[-6pt]
\noindent\rule{\textwidth}{0.5pt}\\[7pt]
{\bfseries RAPPORT DE TRAVAUX PRATIQUES}\\[7pt]
\noindent\rule{\textwidth}{0.5pt}\\[-6pt]
\noindent\rule{\textwidth}{1.6pt}

\vspace{1cm}

{\bfseries T\'el\'ecommunications et R\'eseaux, 4\`eme Ann\'ee}

\vspace{2.2cm}

%%% Titre
{\color{blueESP}\Huge D\'ecouverte et administration\\[6pt] d'un syst\`eme Linux \par}

\vspace{2.8cm}

{\small Par :}

\vspace{0.4cm}

{\normalsize TINOGNY Yonni}

\vspace{1.6cm}

{\small Enseignant :}

\vspace{0.4cm}

{\normalsize Mr. RAKOTOARIJAONA Raonirivo}

\vfill

{\itshape Ann\'ee Universitaire 2025 - 2026}

\vspace{0.5cm}
\end{titlepage}

\tableofcontents
\newpage

% ===================== INTRODUCTION =====================
\section{Introduction}

Ce TP porte sur la d\'ecouverte et l'administration d'un syst\`eme Linux, dans le cadre du parcours
T\'el\'ecommunications et R\'eseaux. Linux constitue l'environnement de r\'ef\'erence dans ce domaine, que ce
soit pour les routeurs, les serveurs, les \'equipements op\'erateurs ou les outils de virtualisation comme GNS3.
La ma\^itrise de cet environnement est donc une comp\'etence de base pour la suite du cursus et pour
l'insertion professionnelle.

Le travail s'organise en quatre parties principales : l'installation d'un syst\`eme Linux en cohabitation
avec Windows, la gestion des utilisateurs et des groupes, l'installation d'applications utiles au cursus, et
enfin une s\'erie de manipulations correspondant \`a un niveau d'administration interm\'ediaire (processus,
r\'eseau, journaux, automatisation par script shell).

Ce rapport pr\'esente les d\'emarches suivies, les commandes utilis\'ees et leur justification, ainsi que les
r\'esultats obtenus \`a chaque \'etape.

% ===================== OBJECTIFS =====================
\section{Objectifs p\'edagogiques}

\`A l'issue de ce TP, l'\'etudiant doit \^etre capable d'installer un syst\`eme Linux en dual-boot avec Windows,
en ma\^itrisant le partitionnement de disque et l'amor\c{c}age via GRUB. Il doit \'egalement savoir cr\'eer et
administrer des comptes utilisateurs et des groupes, et comprendre le mod\`ele de permissions Unix.

L'installation, la configuration et la v\'erification d'un ensemble d'applications professionnelles utiles \`a un
\'etudiant en T\'el\'ecommunications et R\'eseaux font aussi partie des comp\'etences vis\'ees, tout comme la
manipulation des commandes fondamentales et interm\'ediaires de l'administration syst\`eme (processus,
services, paquets, r\'eseau, journaux). Le TP demande enfin d'automatiser une t\^ache r\'ecurrente \`a l'aide
d'un script shell Bash, afin de d\'evelopper l'autonomie et la rigueur m\'ethodologique n\'ecessaires \`a la
poursuite d'\'etudes et \`a l'insertion professionnelle.

% ===================== PRÉREQUIS ET MATÉRIEL =====================
\section{Pr\'erequis et environnement de travail}

Ce TP suppose des notions de base sur l'architecture d'un ordinateur (disque dur, BIOS/UEFI,
partitions) et sur le r\^ole du noyau et du syst\`eme de fichiers dans un syst\`eme d'exploitation. Une
utilisation \'el\'ementaire d'un terminal est un atout, sans \^etre obligatoire.

Le mat\'eriel demand\'e par le sujet comprend un ordinateur sous Windows 10 ou 11 disposant d'au moins
60 Go d'espace libre, une cl\'e USB de 8 Go minimum, une image ISO d'Ubuntu 24.04 LTS, un logiciel de
cr\'eation de cl\'e USB bootable et un acc\`es Internet.

Le poste effectivement utilis\'e est un HP ProBook 450 G5 \'equip\'e de deux disques :

\begin{itemize}[noitemsep]
    \item un SSD de 256 Go (\texttt{sdb}) portant d\'ej\`a Windows 11 et Debian 13 en dual-boot ;
    \item un disque dur de 500 Go (\texttt{sda}, WDC WD5000LPLX) d\'ecoup\'e en deux volumes NTFS de donn\'ees.
\end{itemize}

La machine d\'emarre en mode UEFI. Le syst\`eme install\'e est Ubuntu 24.04 LTS (nom de code \emph{noble}),
\`a partir de l'image \texttt{ubuntu-24.04-desktop-amd64\_3.iso}.

% ===================== PARTIE 1 =====================
\section{Installation d'Ubuntu 24.04 en multi-boot}

\subsection{Adaptation du sujet au poste utilis\'e}

Le sujet part d'un poste sous Windows seul et demande de r\'eduire la partition C: pour y installer Linux.
Ce sc\'enario ne s'appliquait pas ici, puisque le SSD contenait d\'ej\`a deux syst\`emes. R\'eduire encore
les partitions du SSD aurait \'et\'e risqu\'e et aurait laiss\'e peu de marge \`a Windows.

L'objectif a donc \'et\'e adapt\'e : installer Ubuntu comme troisi\`eme syst\`eme, en r\'ecup\'erant de
l'espace sur le disque dur secondaire. Le principe du TP reste le m\^eme, c'est-\`a-dire lib\'erer un espace
non allou\'e, cr\'eer les partitions \`a la main et int\'egrer le nouveau syst\`eme au menu d'amor\c{c}age.
Seul le disque cible change.

\subsection{Pr\'eparation du disque et de la cl\'e USB}

Le partitionnement a \'et\'e r\'ealis\'e avec GParted depuis Debian, et non depuis la Gestion des disques de
Windows. Travailler hors ligne sur les partitions vis\'ees \'evite les probl\`emes li\'es au d\'emarrage
rapide de Windows, qui laisse le syst\`eme de fichiers NTFS dans un \'etat incoh\'erent. Les deux volumes
NTFS du disque dur ont \'et\'e r\'eduits pour lib\'erer environ 40 Go non allou\'es en fin de disque, sans
toucher aux donn\'ees existantes.

\begin{figure}[H]
    \centering
    \includegraphics[width=0.8\textwidth]{gparted_partition_hdd.jpeg}
    \caption{Lib\'eration d'un espace non allou\'e sur le disque dur avec GParted}
\end{figure}

La cl\'e USB bootable a \'et\'e pr\'epar\'ee avec Ventoy plut\^ot qu'avec Rufus. Ventoy installe un
chargeur sur la cl\'e et permet ensuite d'y copier plusieurs images ISO, que l'on choisit au d\'emarrage.
C'est plus souple que Rufus, qui r\'e\'ecrit la cl\'e enti\`ere pour une seule image. L'image d'Ubuntu a
\'et\'e s\'electionn\'ee dans le menu Ventoy.

\begin{figure}[H]
    \centering
    \includegraphics[width=0.8\textwidth]{ventoy_choix_iso.jpeg}
    \caption{S\'election de l'image ISO d'Ubuntu depuis le menu Ventoy}
\end{figure}

Sur ce mod\`ele de HP, les touches d'acc\`es ne sont pas celles indiqu\'ees dans le sujet. F10 ouvre le
BIOS/UEFI et F9 ouvre directement le menu de d\'emarrage. C'est depuis ce menu que la cl\'e USB a \'et\'e
choisie. On y voit d\'ej\`a les entr\'ees UEFI des syst\`emes pr\'esents sur la machine.

\begin{figure}[H]
    \centering
    \includegraphics[width=0.7\textwidth]{boot_menu_f9.jpeg}
    \caption{Menu de d\'emarrage du poste, accessible par la touche F9}
\end{figure}

\subsection{Installation et partitionnement manuel}

L'installateur a \'et\'e lanc\'e en mode interactif. L'option d'installation des pilotes propri\'etaires a
\'et\'e coch\'ee, ce qui est n\'ecessaire sur ce poste pour la carte graphique et surtout pour la carte
Wi-Fi, utilis\'ee ensuite dans toute la partie r\'eseau.

Le sujet demande de contr\^oler pr\'ecis\'ement le d\'ecoupage du disque. L'option de partitionnement
manuel a donc \'et\'e retenue, plut\^ot que l'installation automatique \`a c\^ot\'e des syst\`emes
d\'etect\'es. Quatre partitions ont \'et\'e cr\'e\'ees sur l'espace lib\'er\'e du disque \texttt{sda} :

\begin{itemize}[noitemsep]
    \item \texttt{/} en ext4, 28 Go, pour le syst\`eme ;
    \item \texttt{/boot/efi} en FAT32, 1,13 Go, partition syst\`eme EFI propre \`a Ubuntu ;
    \item \texttt{swap}, 4 Go, dimensionn\'e sans hibernation ;
    \item \texttt{/home} en ext4, 9,82 Go, pour les donn\'ees utilisateur.
\end{itemize}

S\'eparer \texttt{/home} de la racine permet de r\'einstaller le syst\`eme plus tard sans perdre les
fichiers personnels. Le SSD portant Windows et Debian n'a pas \'et\'e touch\'e. Le chargeur d'amor\c{c}age a
\'et\'e install\'e sur \texttt{sda}, le disque qui porte la nouvelle partition EFI.

\begin{figure}[H]
    \centering
    \includegraphics[width=0.85\textwidth]{partitionnement_manuel_ubuntu.jpeg}
    \caption{Partitionnement manuel d'Ubuntu : racine, /boot/efi, swap et /home}
\end{figure}

\subsection{V\'erification de l'amor\c{c}age}

La v\'erification a \'et\'e faite de deux fa\c{c}ons. D'abord avec Grub Customizer depuis Debian, qui liste
les entr\'ees du menu et confirme la pr\'esence des trois syst\`emes.

\begin{figure}[H]
    \centering
    \includegraphics[width=0.85\textwidth]{grub_customizer_menu.jpeg}
    \caption{Entr\'ees du menu d'amor\c{c}age list\'ees par Grub Customizer}
\end{figure}

Ensuite au red\'emarrage, sur le menu GRUB lui-m\^eme. Debian, Ubuntu 24.04 LTS sur \texttt{/dev/sda4} et
Windows Boot Manager sur \texttt{/dev/sdb1} y figurent bien. Les trois syst\`emes ont \'et\'e d\'emarr\'es
l'un apr\`es l'autre pour s'assurer qu'aucun n'avait \'et\'e cass\'e par l'installation.

\begin{figure}[H]
    \centering
    \includegraphics[width=0.85\textwidth]{grub.jpg}
    \caption{Menu GRUB au d\'emarrage, avec les trois syst\`emes install\'es}
\end{figure}

La d\'etection des autres syst\`emes est assur\'ee par \texttt{os-prober}, appel\'e automatiquement lors de
la r\'eg\'en\'eration du fichier de configuration de GRUB. Sa sortie est visible pendant les mises \`a jour
de noyau : elle indique le chemin du chargeur de Windows et la partition de Debian.

\begin{figure}[H]
    \centering
    \includegraphics[width=\textwidth]{update_grub_osprober.png}
    \caption{D\'etection de Windows et de Debian par os-prober lors de la mise \`a jour de GRUB}
\end{figure}

\paragraph{R\'esultats obtenus.}
L'espace n\'ecessaire a \'et\'e lib\'er\'e sur le disque dur sans perte de donn\'ees. Ubuntu 24.04 est
install\'e sur \texttt{/dev/sda4}, aux c\^ot\'es de Windows 11 et Debian 13 rest\'es sur le SSD. Le triple
boot fonctionne et la commande \texttt{sudo update-grub} n'a pas eu besoin d'\^etre relanc\'ee \`a la main,
puisque \texttt{os-prober} avait d\'ej\`a r\'ecup\'er\'e les deux autres syst\`emes.

% ===================== PARTIE 2 =====================
\section{Cr\'eation et gestion des utilisateurs}

Le mod\`ele multi-utilisateur est au c\oe{}ur de l'administration Linux. Cette partie porte sur la cr\'eation de
comptes, la gestion des groupes et le syst\`eme de permissions Unix.

\subsection{Cr\'eation et configuration d'un utilisateur}

Un premier utilisateur a \'et\'e cr\'e\'e avec un r\'epertoire personnel et Bash comme shell par d\'efaut :

\begin{lstlisting}[language=bash]
sudo useradd -m -s /bin/bash etudiant1
sudo passwd etudiant1
sudo usermod -aG sudo etudiant1
id etudiant1
cat /etc/passwd | grep etudiant1
\end{lstlisting}

L'option \texttt{-m} cr\'ee le r\'epertoire personnel, et \texttt{-s /bin/bash} fixe le shell de connexion.
Sans \texttt{-m}, le compte existe mais l'utilisateur se retrouve sans dossier de travail. L'option
\texttt{-aG} de \texttt{usermod} ajoute le compte au groupe \texttt{sudo} sans effacer ses autres groupes ;
oublier le \texttt{-a} remplacerait toute la liste, ce qui est une erreur classique.

La v\'erification se fait avec \texttt{id}, qui affiche les identifiants num\'eriques et les groupes, et avec
la ligne correspondante de \texttt{/etc/passwd}, qui donne le r\'epertoire personnel et le shell.

\begin{figure}[H]
    \centering
    \includegraphics[width=\textwidth]{creation_etudiant1.png}
    \caption{V\'erification du compte etudiant1 avec \texttt{id} et \texttt{/etc/passwd}}
\end{figure}

\paragraph{R\'esultats obtenus.}
La commande \texttt{id etudiant1} renvoie \texttt{uid=1001(etudiant1) gid=1001(etudiant1)
groups=1001(etudiant1),27(sudo)}. L'appartenance au groupe 27, qui est le groupe \texttt{sudo}, confirme que
les droits d'administration ont bien \'et\'e accord\'es. La ligne de \texttt{/etc/passwd} est
\texttt{etudiant1:x:1001:1001::/home/etudiant1:/bin/bash}. Le champ mot de passe contient un \texttt{x},
car l'empreinte du mot de passe est stock\'ee dans \texttt{/etc/shadow}, lisible par root uniquement.

\subsection{Compte \`a usage restreint}

Un second utilisateur, \og invite \fg{}, a \'et\'e cr\'e\'e sans droits d'administration. Il sert \`a
illustrer la gestion du cycle de vie d'un compte : d\'efinition d'une date d'expiration, verrouillage puis
d\'everrouillage.

\begin{lstlisting}[language=bash]
sudo useradd -m -s /bin/bash invite
sudo passwd invite
sudo chage -E 2026-12-01 invite
sudo passwd -l invite
sudo passwd -u invite
\end{lstlisting}

\begin{figure}[H]
    \centering
    \includegraphics[width=\textwidth]{compte_invite.png}
    \caption{Cr\'eation du compte invite, expiration puis verrouillage et d\'everrouillage}
\end{figure}

\paragraph{R\'esultats obtenus.}
Au moment de d\'efinir le mot de passe, le module de qualit\'e a renvoy\'e l'avertissement
\texttt{BAD PASSWORD: The password fails the dictionary check}. Comme la commande \'etait lanc\'ee avec
\texttt{sudo}, le mot de passe a tout de m\^eme \'et\'e accept\'e : root peut passer outre cette v\'erification,
alors qu'un utilisateur ordinaire aurait \'et\'e bloqu\'e.

La commande \texttt{chage -E 2026-12-01} fixe la date au-del\`a de laquelle le compte ne pourra plus se
connecter, ce qui est utile pour un acc\`es temporaire. Les commandes \texttt{passwd -l} et \texttt{passwd -u}
r\'epondent toutes les deux \texttt{password changed}. Le verrouillage n'efface pas le mot de passe : il place
simplement un point d'exclamation devant son empreinte dans \texttt{/etc/shadow}, ce qui rend toute
comparaison impossible. Le d\'everrouillage retire ce caract\`ere et restaure l'acc\`es.

\subsection{Gestion des groupes}

Un groupe de travail a ensuite \'et\'e cr\'e\'e, puis renseign\'e avec l'utilisateur concern\'e :

\begin{lstlisting}[language=bash]
sudo groupadd reseaux4a
sudo usermod -aG reseaux4a etudiant1
getent group reseaux4a
\end{lstlisting}

\texttt{getent} interroge les bases du syst\`eme (fichiers locaux, LDAP, NIS) selon l'ordre d\'efini dans
\texttt{/etc/nsswitch.conf}. Elle est donc pr\'ef\'erable \`a un simple \texttt{cat /etc/group}, qui ne verrait
que les groupes locaux.

\begin{figure}[H]
    \centering
    \includegraphics[width=\textwidth]{gestion_groupe.png}
    \caption{Cr\'eation du groupe reseaux4a et v\'erification de ses membres}
\end{figure}

\paragraph{R\'esultats obtenus.}
La sortie \texttt{reseaux4a:x:1003:etudiant1} montre que le groupe porte le GID 1003 et compte un seul
membre secondaire. Le groupe primaire d'un utilisateur n'appara\^it pas dans cette liste, ce qui explique
que seul \texttt{etudiant1} y figure.

\subsection{Permissions et propri\'et\'es des fichiers}

Le syst\`eme de permissions Unix repose sur trois cat\'egories (propri\'etaire, groupe, autres) et trois droits
(lecture r, \'ecriture w, ex\'ecution x). Un dossier de travail partag\'e a \'et\'e cr\'e\'e pour illustrer ce
m\'ecanisme :

\begin{lstlisting}[language=bash]
sudo mkdir -p /projets/reseaux
sudo chown etudiant1:reseaux4a /projets/reseaux
sudo chmod 770 /projets/reseaux
ls -ld /projets/reseaux
sudo chmod g+s /projets/reseaux
ls -ld /projets/reseaux
\end{lstlisting}

Le mode \texttt{770} accorde lecture, \'ecriture et ex\'ecution au propri\'etaire (7) et au groupe (7), et
aucun droit aux autres (0). Sur un r\'epertoire, le droit d'ex\'ecution correspond en fait au droit de
traverser le dossier, donc d'acc\'eder \`a son contenu. Sans lui, les droits de lecture ne servent \`a rien.

Le bit setgid, pos\'e par \texttt{g+s}, change le comportement par d\'efaut : tout fichier cr\'e\'e dans le
dossier appartient au groupe du dossier parent et non au groupe primaire de son cr\'eateur. C'est ce qui rend
le partage coh\'erent dans la dur\'ee. Sa pr\'esence se lit dans la sortie de \texttt{ls -ld}, o\`u le
\texttt{x} des droits du groupe est remplac\'e par un \texttt{s}, ce qui donne \texttt{drwxrws---}.

% INSÉRER ICI : capture de "ls -ld /projets/reseaux" avant et après le chmod g+s
% \begin{figure}[H]
%     \centering
%     \includegraphics[width=0.85\textwidth]{permissions_dossier.png}
%     \caption{Droits du dossier /projets/reseaux avant et apr\`es l'activation du bit setgid}
% \end{figure}

% ===================== PARTIE 3 =====================
\section{Installation des applications utiles}

\subsection{Mise \`a jour pr\'ealable du syst\`eme}

Avant toute installation, la liste des paquets et le syst\`eme ont \'et\'e mis \`a jour :

\begin{lstlisting}[language=bash]
sudo apt update && sudo apt upgrade -y
\end{lstlisting}

\texttt{apt update} ne fait que rafra\^ichir les index des d\'ep\^ots ; c'est \texttt{apt upgrade} qui
installe r\'eellement les nouvelles versions. Les deux commandes sont encha\^in\'ees avec \texttt{\&\&} pour
que la seconde ne s'ex\'ecute que si la premi\`ere a r\'eussi.

\begin{figure}[H]
    \centering
    \includegraphics[width=\textwidth]{maj_systeme.png}
    \caption{Mise \`a jour du syst\`eme apr\`es l'installation}
\end{figure}

Sur une installation neuve, 51,3 Mo d'index ont \'et\'e t\'el\'echarg\'es et 658 paquets \'etaient \`a mettre
\`a jour. La mise \`a niveau a aussi install\'e un noyau plus r\'ecent (s\'erie HWE), ce qui a d\'eclench\'e la
reg\'en\'eration de l'initramfs et du menu GRUB vue plus haut. Le nettoyage a ensuite \'et\'e fait avec :

\begin{lstlisting}[language=bash]
sudo apt autoremove && sudo apt autoclean
\end{lstlisting}

\subsection{Applications install\'ees}

Les applications demand\'ees par le sujet ont ensuite \'et\'e install\'ees. Chacune r\'epond \`a un besoin
pr\'ecis de la fili\`ere :

\begin{itemize}[noitemsep]
    \item \textbf{VLC} : lecture de flux audio et vid\'eo, tests de streaming RTP, RTSP et UDP ;
    \item \textbf{Evince} : lecture des polycopi\'es et des rapports au format PDF ;
    \item \textbf{TeX Live et TeXstudio} : r\'edaction des rapports scientifiques en LaTeX ;
    \item \textbf{Visual Studio Code} : \'edition de scripts shell, de code Python ou C et de fichiers de configuration ;
    \item \textbf{GNS3} : simulation de topologies r\'eseau et de protocoles de t\'el\'ecommunication ;
    \item \textbf{netsniff-ng} : capture de trafic bas niveau, en compl\'ement de Wireshark ;
    \item \textbf{Wireshark} : analyse fine des protocoles r\'eseau ;
    \item \textbf{OpenSSH Server} : administration \`a distance, utilis\'ee dans les TP de supervision ;
    \item \textbf{Git} : gestion de versions des projets et des scripts ;
    \item \textbf{build-essential} : cha\^ine de compilation C et C++.
\end{itemize}

\begin{lstlisting}[language=bash]
sudo apt install vlc
sudo apt install evince
sudo apt install texlive-full texstudio
sudo apt install wireshark
sudo apt install netsniff-ng
sudo apt install openssh-server
sudo apt install git
sudo apt install build-essential
sudo apt install htop
\end{lstlisting}

\subsection{Cas particuliers rencontr\'es}

Trois installations n'ont pas suivi exactement la proc\'edure du sujet.

\paragraph{Visual Studio Code.}
Le sujet propose \texttt{sudo snap install code --classic}. Le paquet Debian officiel de Microsoft a \'et\'e
pr\'ef\'er\'e, car il ajoute le d\'ep\^ot \texttt{packages.microsoft.com} au syst\`eme et permet donc de
suivre les mises \`a jour avec \texttt{apt}, sans passer par un second gestionnaire de paquets.

\begin{lstlisting}[language=bash]
sudo apt install ./code_1.131.0-1785237861_amd64.deb
\end{lstlisting}

\begin{figure}[H]
    \centering
    \includegraphics[width=\textwidth]{install_vscode.png}
    \caption{Installation de Visual Studio Code \`a partir du paquet .deb officiel}
\end{figure}

La capture montre aussi un message r\'ep\'et\'e \texttt{Waiting for cache lock}. Une autre instance d'apt
tournait encore dans un second onglet du terminal. Le verrou \texttt{/var/lib/dpkg/lock-frontend} garantit
qu'un seul processus modifie la base des paquets \`a la fois. Il a suffi d'attendre la fin de l'op\'eration
pr\'ec\'edente.

\paragraph{GNS3.}
Les paquets \texttt{gns3-gui} et \texttt{gns3-server} ne sont pas pr\'esents dans les d\'ep\^ots standards
d'Ubuntu 24.04. Le d\'ep\^ot PPA du projet a donc \'et\'e ajout\'e avant l'installation :

\begin{lstlisting}[language=bash]
sudo add-apt-repository ppa:gns3/ppa
sudo apt install gns3-gui gns3-server
\end{lstlisting}

\begin{figure}[H]
    \centering
    \includegraphics[width=\textwidth]{gns3_ppa.png}
    \caption{Ajout du d\'ep\^ot PPA de GNS3 puis installation des paquets}
\end{figure}

Pendant l'installation, le paquet \texttt{ubridge} demande si les utilisateurs non privil\'egi\'es doivent
pouvoir lancer GNS3. La r\'eponse affirmative a \'et\'e retenue : les membres des groupes \texttt{sudo} et
\texttt{admin} sont alors ajout\'es au groupe \texttt{ubridge}, qui re\c{c}oit les capacit\'es de capture. Cela
\'evite de faire tourner tout GNS3 en tant que root, donc limite la surface d'exposition.

\paragraph{netsniff-ng.}
L'installation de ce paquet entra\^ine la suppression de \texttt{systemd-timesyncd}, remplac\'e par
\texttt{ntpsec}. Les deux services assurent la synchronisation horaire et ne peuvent pas coexister. Ce
changement est sans cons\'equence ici, mais il montre l'int\'er\^et de lire le r\'ecapitulatif propos\'e par
apt avant de valider, en particulier la ligne \og The following packages will be REMOVED \fg{}.

\subsection{V\'erification des installations}

Le bon fonctionnement a \'et\'e contr\^ol\'e soit par le lancement de l'interface graphique, soit en ligne de
commande :

\begin{lstlisting}[language=bash]
apt list --installed | grep vlc
code --version
vlc --version
wireshark --version
git --version
gcc --version
\end{lstlisting}

\begin{figure}[H]
    \centering
    \includegraphics[width=\textwidth]{verif_vlc.png}
    \caption{V\'erification de l'installation de VLC et de ses greffons}
\end{figure}

\paragraph{R\'esultats obtenus.}
Toutes les applications demand\'ees ont \'et\'e install\'ees. La v\'erification de VLC renvoie la version
3.0.20-3build6 pour le paquet principal et pour l'ensemble de ses greffons, tous marqu\'es
\texttt{[installed]} ou \texttt{[installed,automatic]}. La mention \texttt{automatic} indique les paquets
tir\'es comme d\'ependances, qui seront retir\'es automatiquement par \texttt{apt autoremove} si VLC est
d\'esinstall\'e. L'avertissement \og apt does not have a stable CLI interface \fg{} rappelle que
\texttt{apt} est pr\'evu pour un usage interactif ; dans un script, il vaut mieux utiliser
\texttt{apt-get} ou \texttt{dpkg-query}.

% ===================== PARTIE 4 =====================
\section{Manipulations de niveau interm\'ediaire}

Cette partie regroupe les comp\'etences qui distinguent un utilisateur d\'ebutant d'un utilisateur
interm\'ediaire, avec un accent particulier sur l'administration r\'eseau.

\subsection{Gestion des processus}

\begin{lstlisting}[language=bash]
ps aux
ps aux | grep obsidian
top
htop
kill -9 <PID>
\end{lstlisting}

\texttt{ps aux} donne une photographie des processus \`a un instant donn\'e, alors que \texttt{top} et
\texttt{htop} rafra\^ichissent l'affichage en continu. En pratique, on filtre presque toujours la sortie de
\texttt{ps} avec \texttt{grep} pour retrouver un processus pr\'ecis.

\begin{figure}[H]
    \centering
    \includegraphics[width=\textwidth]{processus_kill.png}
    \caption{Recherche d'un processus avec \texttt{ps aux | grep} puis arr\^et forc\'e avec \texttt{kill -9}}
\end{figure}

\paragraph{R\'esultats obtenus.}
L'application test\'ee appara\^it sous la forme de plusieurs processus : un processus principal et des
processus fils charg\'es du rendu et du r\'eseau. Le PID du processus principal (18030) a \'et\'e transmis \`a
\texttt{kill -9}, ce qui a fait dispara\^itre toute l'arborescence. La derni\`ere ligne de la sortie
correspond au \texttt{grep} lui-m\^eme, qui appara\^it dans sa propre recherche.

Le signal 9 (\texttt{SIGKILL}) ne peut pas \^etre intercept\'e par le programme, qui est arr\^et\'e sans
pouvoir enregistrer son \'etat. Il vaut donc mieux essayer d'abord \texttt{kill <PID>}, qui envoie
\texttt{SIGTERM} et laisse le processus se fermer proprement.

\subsection{Gestion des services}

Le service SSH a servi d'exemple, car il est directement utile pour la suite du cursus.

\begin{lstlisting}[language=bash]
systemctl status ssh
sudo systemctl start ssh
sudo systemctl enable ssh
\end{lstlisting}

\begin{figure}[H]
    \centering
    \includegraphics[width=\textwidth]{statut_ssh.png}
    \caption{\'Etat du service SSH avant et apr\`es d\'emarrage}
\end{figure}

\paragraph{R\'esultats obtenus.}
Avant d\'emarrage, le service est \texttt{inactive (dead)} et marqu\'e \texttt{disabled}. Apr\`es
\texttt{systemctl start}, il passe \`a \texttt{active (running)} avec le PID 20047, et les journaux
confirment l'\'ecoute sur le port 22 en IPv4 et en IPv6.

La commande \texttt{systemctl enable ssh} a ensuite cr\'e\'e les liens symboliques dans
\texttt{/etc/systemd/system/multi-user.target.wants/}, ce qui rend le d\'emarrage automatique au prochain
amor\c{c}age. Le champ \texttt{Loaded} passe alors de \texttt{disabled} \`a \texttt{enabled}. La distinction
est importante : \texttt{start} agit sur la session en cours, \texttt{enable} agit sur les d\'emarrages
suivants.

On remarque aussi la ligne \texttt{TriggeredBy: ssh.socket}. Ubuntu active le service \`a la demande via un
socket systemd, ce qui explique que le service puisse repasser \`a \texttt{inactive} tant qu'aucune connexion
n'arrive.

\subsection{Gestion des paquets}

\begin{lstlisting}[language=bash]
apt list --installed | grep vlc
sudo apt remove <paquet>
sudo apt autoremove
dpkg -l
\end{lstlisting}

Ces commandes permettent de v\'erifier qu'un paquet est install\'e, de le d\'esinstaller, de nettoyer les
d\'ependances devenues inutiles, et de lister l'ensemble des paquets du syst\`eme. \texttt{apt remove} laisse
les fichiers de configuration en place ; \texttt{apt purge} les supprime \'egalement.

\subsection{Administration r\'eseau}

\subsubsection{\'Etat des interfaces}

\begin{lstlisting}[language=bash]
ip a
ip route
ss -tuln
ping -c 4 8.8.8.8
traceroute www.univ.mg
\end{lstlisting}

\begin{figure}[H]
    \centering
    \includegraphics[width=\textwidth]{ip_a.png}
    \caption{Interfaces r\'eseau et adresses IP du poste}
\end{figure}

\paragraph{R\'esultats obtenus.}
Le poste pr\'esente quatre interfaces. \texttt{lo} est l'interface de bouclage en 127.0.0.1.
\texttt{enp1s0} est l'interface Ethernet, marqu\'ee \texttt{NO-CARRIER} et \texttt{state DOWN} car aucun
c\^able n'est branch\'e. \texttt{wlp2s0} est l'interface Wi-Fi, qui porte l'adresse 10.251.134.190/24 obtenue
en DHCP, comme l'indiquent le mot-cl\'e \texttt{dynamic} et le bail de 3550 secondes. \texttt{virbr0} est
un pont virtuel en 192.168.122.1/24 cr\'e\'e par libvirt, utilis\'e par les machines virtuelles.

\subsubsection{Configuration d'une adresse IP statique avec Netplan}

Netplan lit les fichiers YAML de \texttt{/etc/netplan/} et g\'en\`ere la configuration du gestionnaire de
r\'eseau sous-jacent. Sur une version bureau d'Ubuntu, ce gestionnaire est NetworkManager, ce qui impose de
pr\'eciser \texttt{renderer: NetworkManager}. Le fichier suivant a \'et\'e cr\'e\'e :

\begin{lstlisting}
network:
  version: 2
  renderer: NetworkManager
  wifis:
    wlp2s0:
      dhcp4: no
      access-points:
        "Airbox-4G":
          password: "********"
      addresses:
        - 10.251.134.191/24
      routes:
        - to: default
          via: 10.251.134.118
      nameservers:
        addresses: [8.8.8.8, 1.1.1.1]
\end{lstlisting}

\begin{figure}[H]
    \centering
    \includegraphics[width=\textwidth]{netplan_yaml.png}
    \caption{Fichier \texttt{/etc/netplan/99-config-statique.yaml} (mot de passe masqu\'e)}
\end{figure}

L'adresse choisie est voisine de celle attribu\'ee par le DHCP, mais volontairement diff\'erente pour
\'eviter un conflit. Le fichier est num\'erot\'e 99 pour qu'il soit lu en dernier et prenne le pas sur la
configuration par d\'efaut. Il a ensuite \'et\'e trait\'e ainsi :

\begin{lstlisting}[language=bash]
sudo chmod 600 /etc/netplan/99-config-statique.yaml
sudo netplan generate
sudo netplan try
\end{lstlisting}

Le \texttt{chmod 600} r\'epond \`a un avertissement de Netplan : le fichier contient la cl\'e du r\'eseau
Wi-Fi en clair et ne doit \^etre lisible que par root. \texttt{netplan generate} valide la syntaxe et produit
les fichiers du gestionnaire de r\'eseau. \texttt{netplan try} applique la configuration en laissant un
compte \`a rebours ; sans confirmation, l'ancienne configuration est restaur\'ee. C'est une pr\'ecaution
utile quand on travaille \`a distance, o\`u une erreur couperait l'acc\`es \`a la machine.

\paragraph{Premi\`ere tentative, en \'echec.}
L'activation du profil g\'en\'er\'e a \'echou\'e avec le message \texttt{Connection activation failed: The
Wi-Fi network could not be found}.

\begin{figure}[H]
    \centering
    \includegraphics[width=\textwidth]{netplan_echec.png}
    \caption{\'Echec de l'activation du profil Netplan et adresse toujours obtenue en DHCP}
\end{figure}

Le nom du r\'eseau inscrit dans le fichier ne correspondait pas au point d'acc\`es r\'eellement port\'ee. La
v\'erification avec \texttt{ip a show wlp2s0} le confirme : l'interface conservait l'adresse 10.251.134.190
en mode \texttt{dynamic}, donc toujours fournie par le DHCP. La configuration \'etait syntaxiquement
correcte, mais inapplicable dans ce contexte.

\paragraph{Seconde tentative, r\'eussie.}
La manipulation a \'et\'e reprise sur un r\'eseau disponible. Le fichier a
\'et\'e adapt\'e au SSID et au plan d'adressage de ce r\'eseau, puis le profil correspondant a \'et\'e
activ\'e :

\begin{lstlisting}[language=bash]
nmcli connection show --active
sudo nmcli connection up "netplan-wlp2s0-3A-XL"
ip a show wlp2s0
\end{lstlisting}

\begin{figure}[H]
    \centering
    \includegraphics[width=\textwidth]{ip_statique.png}
    \caption{Activation du profil Netplan et adresse statique effectivement appliqu\'ee}
\end{figure}

\paragraph{R\'esultats obtenus.}
Le message \texttt{Connection successfully activated} confirme l'activation. L'interface porte ensuite
l'adresse 10.151.30.191/24. Le mot-cl\'e \texttt{dynamic} a disparu de la sortie de \texttt{ip a}, et le bail
DHCP n'appara\^it plus : l'adresse est bien fixe. La mention \texttt{noprefixroute} indique que la route du
r\'eseau local n'est pas ajout\'ee automatiquement, la table de routage \'etant construite \`a partir de la
section \texttt{routes} du fichier.

\subsubsection{Pare-feu}

Un pare-feu de base a \'et\'e mis en place avec UFW, en n'ouvrant que le port n\'ecessaire \`a
l'administration \`a distance :

\begin{lstlisting}[language=bash]
sudo ufw allow 22/tcp
sudo ufw enable
sudo ufw status verbose
\end{lstlisting}

L'ordre des commandes a son importance. Autoriser le port 22 avant d'activer le pare-feu \'evite de couper
une session SSH en cours, puisque la politique par d\'efaut d'UFW bloque tout le trafic entrant.

% INSÉRER ICI : capture de "sudo ufw status verbose"
% \begin{figure}[H]
%     \centering
%     \includegraphics[width=0.85\textwidth]{ufw_status.png}
%     \caption{R\`egles actives du pare-feu UFW}
% \end{figure}

\subsubsection{Capture de trafic}

\begin{lstlisting}[language=bash]
sudo tcpdump -i wlp2s0 -n
sudo netsniff-ng -i wlp2s0 -o capture.pcap
\end{lstlisting}

Les exemples du sujet utilisent \texttt{eth0}. Ce nom n'existe plus sur Ubuntu 24.04, qui applique la
convention de nommage pr\'evisible des interfaces : \texttt{enp1s0} pour l'Ethernet et \texttt{wlp2s0} pour le
Wi-Fi. Les noms d\'ependent de l'emplacement mat\'eriel du p\'eriph\'erique, ce qui \'evite qu'ils changent
d'un d\'emarrage \`a l'autre. Les commandes ont donc \'et\'e adapt\'ees \`a l'interface r\'eellement active.

L'option \texttt{-n} de \texttt{tcpdump} d\'esactive la r\'esolution DNS des adresses. L'affichage est plus
rapide et, surtout, la capture ne g\'en\`ere pas elle-m\^eme de trafic suppl\'ementaire.

\subsection{Simulation sous GNS3}

GNS3 a \'et\'e install\'e depuis le d\'ep\^ot du projet et se lance correctement en version 2.2.60. La
fen\^etre des pr\'ef\'erences donne acc\`es aux diff\'erentes familles d'\'equipements \'emulables :
routeurs Dynamips, machines Qemu, conteneurs Docker et commutateurs int\'egr\'es.

\begin{figure}[H]
    \centering
    \includegraphics[width=0.85\textwidth]{gns3_lance.png}
    \caption{GNS3 2.2.60 lanc\'e, fen\^etre des mod\`eles de routeurs IOS}
\end{figure}

\subsubsection{Topologie r\'ealis\'ee}

La topologie demand\'ee par le sujet comporte deux routeurs, un commutateur et deux PC. Les routeurs sont
\'emul\'es par Dynamips \`a partir d'une image IOS, le commutateur est le mod\`ele Ethernet int\'egr\'e \`a
GNS3, et les deux postes utilisent VPCS. VPCS n'est pas une machine virtuelle compl\`ete : c'est un
simulateur de station tr\`es l\'eger qui sait seulement recevoir une adresse, r\'epondre au ping et lancer un
traceroute. C'est suffisant ici et bien moins co\^uteux en ressources qu'une VM.

Le montage relie PC1 au commutateur, le commutateur \`a la premi\`ere interface de R1, les deux routeurs
entre eux, puis PC2 directement sur R2. Deux r\'eseaux locaux distincts sont ainsi s\'epar\'es par une
liaison de transit, ce qui oblige \`a faire du routage pour les faire communiquer.

\begin{center}
\begin{tabular}{llll}
\hline
\textbf{\'Equipement} & \textbf{Interface} & \textbf{Adresse IP} & \textbf{R\^ole} \\
\hline
PC1 & e0      & 192.168.1.10/24 & Station du LAN 1 \\
R1  & f0/0    & 192.168.1.1/24  & Passerelle du LAN 1 \\
R1  & f0/1    & 10.0.0.1/30     & Liaison inter-routeurs \\
R2  & f0/1    & 10.0.0.2/30     & Liaison inter-routeurs \\
R2  & f0/0    & 192.168.2.1/24  & Passerelle du LAN 2 \\
PC2 & e0      & 192.168.2.10/24 & Station du LAN 2 \\
\hline
\end{tabular}
\end{center}

Une liaison en /30 a \'et\'e choisie entre les deux routeurs. Ce masque ne laisse que deux adresses
utilisables, ce qui correspond exactement au besoin d'un lien point \`a point et \'evite de gaspiller un
r\'eseau entier.

La configuration des adresses sur les PC se fait directement dans la console VPCS :

\begin{lstlisting}
PC1> ip 192.168.1.10 255.255.255.0 192.168.1.1
PC2> ip 192.168.2.10 255.255.255.0 192.168.2.1
\end{lstlisting}

C\^ot\'e routeurs, chaque interface a \'et\'e adress\'ee puis activ\'ee. Par d\'efaut une interface IOS est
administrativement d\'esactiv\'ee, d'o\`u la commande \texttt{no shutdown} qu'il ne faut pas oublier :

\begin{lstlisting}
R1(config)# interface fastEthernet 0/0
R1(config-if)# ip address 192.168.1.1 255.255.255.0
R1(config-if)# no shutdown
R1(config-if)# exit
R1(config)# ip route 192.168.2.0 255.255.255.0 10.0.0.2
\end{lstlisting}

La derni\`ere ligne est indispensable. Un routeur ne conna\^it au d\'epart que les r\'eseaux directement
connect\'es \`a ses interfaces. Sans route vers le LAN oppos\'e, R1 re\c{c}oit bien les paquets de PC1 mais
ne sait pas o\`u les envoyer et les rejette. Une route statique sym\'etrique a donc \'et\'e ajout\'ee sur R2
vers 192.168.1.0/24.

\paragraph{R\'esultats obtenus.}
La validation s'est faite en trois temps, du plus proche au plus lointain. Le ping de PC1 vers sa passerelle
192.168.1.1 confirme d'abord que le lien local et le commutateur fonctionnent. Le ping de R1 vers 10.0.0.2
v\'erifie ensuite la liaison inter-routeurs. Le ping de PC1 vers PC2, enfin, aboutit \'egalement, ce qui
prouve que les deux routes statiques sont correctes et que le routage entre les deux LAN est
op\'erationnel.

Cette progression est la m\'ethode \`a retenir : tester segment par segment plut\^ot que de lancer
directement le ping de bout en bout. En cas d'\'echec, elle indique imm\'ediatement quel maillon est en
cause.

\subsection{Journaux syst\`emes et supervision}

\begin{lstlisting}[language=bash]
journalctl -xe
journalctl -u ssh
df -h
du -sh /home/*
free -h
\end{lstlisting}

\texttt{journalctl} lit le journal binaire de systemd. L'option \texttt{-u} filtre sur une unit\'e pr\'ecise,
ce qui est le r\'eflexe \`a avoir quand un service refuse de d\'emarrer. Les commandes \texttt{df},
\texttt{du} et \texttt{free} renseignent respectivement sur l'espace disque disponible, la taille des
r\'epertoires utilisateurs et l'occupation de la m\'emoire. L'option \texttt{-h} affiche des unit\'es
lisibles plut\^ot que des blocs.

Sur ce poste, \texttt{df -h} est utile pour surveiller la partition racine, limit\'ee \`a 28 Go, que
l'installation de \texttt{texlive-full} et de GNS3 remplit rapidement.

\subsection{Recherche et traitement de texte}

\begin{lstlisting}[language=bash]
grep -i 'error' /var/log/syslog
sed 's/ancien/nouveau/g' fichier.txt
awk -F: '{print $1}' /etc/passwd
find / -name '*.log' -mtime -1
\end{lstlisting}

\texttt{grep} recherche un motif, ici les lignes contenant \og error \fg{} sans distinction de casse gr\^ace
\`a \texttt{-i}. \texttt{sed} remplace du texte, le \texttt{g} final portant le remplacement sur toutes les
occurrences d'une m\^eme ligne et non seulement sur la premi\`ere. \texttt{awk} d\'ecoupe chaque ligne selon
un s\'eparateur, ici le deux-points, et n'affiche que le premier champ, c'est-\`a-dire la liste des noms
d'utilisateurs. \texttt{find} recherche des fichiers selon des crit\`eres, ici les journaux modifi\'es depuis
moins de 24 heures.

Ces quatre commandes sont la base du d\'epannage sur un \'equipement distant, o\`u l'on ne dispose que d'un
terminal pour fouiller des fichiers de plusieurs milliers de lignes.

\subsection{Archivage et compression}

\begin{lstlisting}[language=bash]
tar -czvf archive.tar.gz /projets/reseaux
tar -xzvf archive.tar.gz
\end{lstlisting}

La premi\`ere commande cr\'ee une archive compress\'ee, la seconde l'extrait. Les options se lisent une \`a
une : \texttt{c} pour cr\'eer, \texttt{x} pour extraire, \texttt{z} pour la compression gzip, \texttt{v} pour
l'affichage d\'etaill\'e et \texttt{f} pour d\'esigner le fichier cible.

\subsection{Automatisation par script shell}

Un script Bash nomm\'e \texttt{surveillance\_reseau.sh} a \'et\'e r\'edig\'e pour enregistrer
p\'eriodiquement l'\'etat de la connectivit\'e :

\begin{lstlisting}[language=bash]
#!/bin/bash
# surveillance_reseau.sh
# Teste la connectivite vers un serveur externe et journalise le resultat.

LOG="$HOME/surveillance_reseau.log"
DATE=$(date '+%Y-%m-%d %H:%M:%S')

if ping -c 1 -W 2 8.8.8.8 &> /dev/null; then
    echo "$DATE - Connexion OK" >> "$LOG"
else
    echo "$DATE - Connexion PERDUE" >> "$LOG"
fi
\end{lstlisting}

Quelques points m\'eritent d'\^etre expliqu\'es. La redirection \texttt{\&> /dev/null} envoie la sortie
standard et la sortie d'erreur vers le p\'eriph\'erique nul : seul le code de retour de \texttt{ping}
int\'eresse le script. L'option \texttt{-c 1} limite l'envoi \`a un seul paquet, et \texttt{-W 2} borne
l'attente \`a deux secondes pour que le script ne reste pas bloqu\'e si le r\'eseau est coup\'e.

Le sujet propose d'\'ecrire dans \texttt{/var/log/surveillance\_reseau.log}. Ce chemin n'est accessible en
\'ecriture qu'\`a root, alors qu'une t\^ache plac\'ee dans la crontab d'un utilisateur s'ex\'ecute avec les
droits de cet utilisateur. Le fichier de log a donc \'et\'e plac\'e dans le r\'epertoire personnel. Pour le
conserver dans \texttt{/var/log}, il faudrait soit passer par la crontab de root, soit cr\'eer le fichier au
pr\'ealable et en donner la propri\'et\'e \`a l'utilisateur.

Le script a ensuite \'et\'e rendu ex\'ecutable, test\'e \`a la main, puis planifi\'e :

\begin{lstlisting}[language=bash]
chmod +x surveillance_reseau.sh
./surveillance_reseau.sh
cat ~/surveillance_reseau.log
crontab -e
\end{lstlisting}

\begin{lstlisting}
# Execution toutes les 5 minutes
*/5 * * * * /home/etudiant1/surveillance_reseau.sh
\end{lstlisting}

Les cinq champs d'une ligne crontab correspondent aux minutes, aux heures, au jour du mois, au mois et au
jour de la semaine. La notation \texttt{*/5} en premi\`ere position d\'eclenche l'ex\'ecution toutes les cinq
minutes. Le chemin du script doit \^etre absolu, car cron s'ex\'ecute dans un environnement minimal o\`u la
variable \texttt{PATH} est tr\`es r\'eduite.

% INSÉRER ICI : capture de l'exécution du script, de la crontab et du fichier de log
% \begin{figure}[H]
%     \centering
%     \includegraphics[width=0.85\textwidth]{execution_script.png}
%     \caption{Ex\'ecution du script et contenu du fichier de log}
% \end{figure}

Cette combinaison de script et de planification reproduit, \`a petite \'echelle, ce que font les sondes de
supervision d\'eploy\'ees chez un op\'erateur : un test simple, r\'ep\'et\'e \`a intervalle fixe, dont la
trace horodat\'ee permet de reconstituer les p\'eriodes d'indisponibilit\'e.

% ===================== ANALYSE CRITIQUE =====================
\section{Analyse critique des difficult\'es rencontr\'ees}

\subsection{Un poste qui ne correspondait pas au sc\'enario du sujet}

La premi\`ere difficult\'e a \'et\'e d'ordre m\'ethodologique. Le sujet d\'ecrit une machine sous Windows
seul, alors que le poste comportait d\'ej\`a deux syst\`emes sur un SSD presque plein. Appliquer la
proc\'edure telle quelle aurait conduit \`a r\'eduire encore les partitions existantes.

Le choix a \'et\'e d'installer Ubuntu sur le disque dur secondaire, apr\`es y avoir lib\'er\'e de l'espace.
Cela demandait de bien identifier les disques, car une confusion entre \texttt{sda} et \texttt{sdb} au moment
d'installer le chargeur d'amor\c{c}age aurait pu rendre les autres syst\`emes inaccessibles. Les captures de
GParted et de l'installateur ont servi de contr\^ole \`a chaque \'etape. Le r\'esultat est un triple boot
fonctionnel.

\subsection{Touches d'acc\`es au BIOS diff\'erentes}

Le sujet indique F12, F2, \'Echap ou Suppr pour atteindre le menu de d\'emarrage. Aucune de ces touches ne
fonctionne sur ce mod\`ele de HP, o\`u il faut appuyer sur F9 pour le menu de d\'emarrage et F10 pour le
BIOS. Le point est mineur, mais il illustre que ces raccourcis d\'ependent du constructeur et qu'il faut
toujours v\'erifier la documentation de la machine.

\subsection{Paquets absents des d\'ep\^ots standards}

GNS3 n'est pas distribu\'e dans les d\'ep\^ots officiels d'Ubuntu 24.04. La commande donn\'ee par le sujet
\'echoue donc telle quelle. L'ajout du PPA du projet a r\'esolu le probl\`eme. Cela pose toutefois une
question de confiance : un PPA est un d\'ep\^ot tiers, et l'ajouter revient \`a autoriser son mainteneur \`a
installer des paquets sur le syst\`eme. Ici le d\'ep\^ot est celui de l'\'editeur du logiciel, ce qui est
acceptable.

Le m\^eme raisonnement a conduit \`a installer Visual Studio Code depuis le paquet officiel de Microsoft
plut\^ot que depuis Snap, afin de garder un seul gestionnaire de paquets \`a suivre.

\subsection{Verrou du gestionnaire de paquets}

Plusieurs installations ont \'echou\'e sur des messages de verrou :
\texttt{Could not get lock /var/lib/dpkg/lock-frontend}. Deux causes distinctes ont \'et\'e rencontr\'ees.
Dans le premier cas, la commande avait \'et\'e lanc\'ee sans \texttt{sudo} et le refus \'etait un simple
probl\`eme de droits. Dans le second, une autre instance d'apt tournait dans un onglet voisin du terminal.

La bonne pratique consiste \`a identifier le processus propri\'etaire du verrou, dont apt affiche le PID,
puis \`a attendre qu'il se termine. Supprimer le fichier de verrou \`a la main est possible mais dangereux :
cela peut laisser la base dpkg dans un \'etat incoh\'erent.

\subsection{Configuration r\'eseau statique sur une interface Wi-Fi}

C'est la difficult\'e la plus instructive de ce TP. Le fichier Netplan \'etait syntaxiquement valide, accept\'e
par \texttt{netplan generate}, et pourtant l'adresse restait fournie par le DHCP. L'activation du profil
retournait \texttt{The Wi-Fi network could not be found}.

L'analyse a montr\'e que le probl\`eme ne venait pas du fichier mais du contexte. Sur une interface Wi-Fi,
Netplan ne configure pas seulement une adresse : il d\'ecrit aussi le point d'acc\`es auquel se connecter.
Si le SSID d\'eclar\'e n'est pas \`a port\'ee, le profil ne peut pas s'activer et NetworkManager conserve la
connexion pr\'ec\'edente. La v\'erification par \texttt{ip a show wlp2s0} a permis de le confirmer, la mention
\texttt{dynamic} indiquant clairement une adresse encore obtenue en DHCP.

La manipulation a \'et\'e reprise sur un autre r\'eseau, avec le SSID et la passerelle correspondants. Le
profil s'est activ\'e et l'adresse statique a bien \'et\'e prise en compte. La le\c{c}on \`a retenir est
qu'une configuration valide n'est pas forc\'ement applicable, et qu'il faut toujours v\'erifier le r\'esultat
sur l'interface plut\^ot que de se fier au seul retour de la commande d'application.

\subsection{Politique de mot de passe}

La cr\'eation du compte \og invite \fg{} a d\'eclench\'e l'avertissement
\texttt{BAD PASSWORD: The password fails the dictionary check}. Le mot de passe a tout de m\^eme \'et\'e
accept\'e, car la commande \'etait lanc\'ee en tant que root. Ce comportement m\'erite d'\^etre connu : un
administrateur peut cr\'eer sans s'en rendre compte des comptes dont le mot de passe ne respecte pas la
politique qu'il impose aux autres utilisateurs.

% ===================== CONCLUSION =====================
\section{Conclusion}

Ce TP a permis de parcourir le cycle complet d'administration d'un poste Linux, depuis l'installation du
syst\`eme jusqu'\`a l'automatisation d'une t\^ache de supervision. L'installation en multi-boot a impos\'e
une rigueur particuli\`ere sur le partitionnement et sur le choix du disque d'amor\c{c}age, o\`u une erreur
aurait pu rendre les autres syst\`emes inaccessibles.

La gestion des utilisateurs et des groupes a mis en \'evidence le mod\`ele de permissions Unix, avec la
distinction entre propri\'etaire, groupe et autres, et le r\^ole du bit setgid dans un dossier partag\'e.
L'installation des applications a pr\'epar\'e l'environnement de travail des prochains TP, tout en montrant
que les commandes d'un \'enonc\'e ne s'appliquent pas toujours telles quelles : d\'ep\^ots manquants, noms
d'interfaces diff\'erents, conflits entre paquets.

Les manipulations de niveau interm\'ediaire ont surtout \'et\'e l'occasion de travailler la m\'ethode. La
configuration Netplan qui ne s'appliquait pas en est le meilleur exemple : la commande ne renvoyait pas
d'erreur de syntaxe, et seule la v\'erification de l'\'etat r\'eel de l'interface a permis de comprendre ce
qui se passait. C'est exactement le r\'eflexe attendu en exploitation r\'eseau, o\`u l'on valide toujours une
configuration par son effet observable et non par le simple retour de la commande.

La simulation sous GNS3 a permis de v\'erifier ces acquis sur une topologie rout\'ee, avec la validation
progressive de la connectivit\'e segment par segment. Il reste \`a approfondir la mise en place d'une
supervision plus compl\`ete que le simple test de connectivit\'e r\'ealis\'e ici, par exemple avec un
outil de collecte comme SNMP. L'ensemble constitue une base solide pour la suite du cursus en
T\'el\'ecommunications et R\'eseaux.

% ===================== BIBLIOGRAPHIE =====================
\section{Ressources bibliographiques et compl\'ementaires}

\begin{itemize}[noitemsep]
    \item Documentation officielle Ubuntu : \url{https://help.ubuntu.com/}
    \item Documentation officielle Debian : \url{https://www.debian.org/doc/}
    \item \textit{The Linux Command Line}, W. Shotts (ouvrage libre)
    \item Documentation Netplan : \url{https://netplan.readthedocs.io/}
    \item Documentation GNS3 : \url{https://docs.gns3.com/}
    \item Manuel des commandes Linux : \texttt{man <commande>}
\end{itemize}

\end{document}
